\section*{Chapitre \ref{ch:iv1:mental-technique} --- Calcul mental}
\markboth{\protect\omTocMark{}{Réponses~: Calcul mental}}{\protect\omTocMark{}{Réponses~: Calcul mental}}

\begin{solution}{iq:iv1:mental-technique:1}
Le produit vaut 9991.

Utiliser la différence de carrés\index{différence de carrés}. Les deux nombres sont à 3 de part et d'autre de 100, donc $97 \times 103 = (100 - 3)(100 + 3) = 100^2 - 3^2 = 9991$.

L'idée fonctionne dès que les deux nombres ont une somme paire et que leur moyenne est facile à élever au carré. Pour $47 \times 53$, la moyenne vaut 50, ce qui donne $2500 - 9 = 2491$. Pour $96 \times 108$, la moyenne vaut 102 et chaque nombre en est à 6, ce qui donne $10\,404 - 36 = 10\,368$.

\iqfollow{Et si la somme est impaire, comme dans $47 \times 54$~?}{On retire d'abord une copie~: $47 \times 54 = 47 \times 53 + 47 = 2491 + 47 = 2538$.}
\iqwrong{Dire «~environ 10\,000~». C'est une bonne vérification de l'ordre de grandeur, mais un recruteur qui demande le produit attend qu'on retranche les 9. La correction est un petit carré unique, elle ne coûte rien.}
\end{solution}

\begin{solution}{iq:iv1:mental-technique:2}
$85^2 = 7225$. La règle~: on multiplie la partie de tête par son successeur, puis on écrit 25.

Pour 85, la partie de tête vaut 8, et $8 \times 9 = 72$, donc la réponse est 72 suivi de 25. Pour le démontrer, on écrit le nombre sous la forme $10a + 5$~:
\[
(10a + 5)^2 = 100a^2 + 100a + 25 = 100\,a(a + 1) + 25.
\]
La règle vaut pour tout $a$, et pas seulement pour un chiffre~: $115^2$ a $a = 11$, et $11 \times 12 = 132$ donne $13\,225$.

\iqfollow{Existe-t-il une règle semblable autour de 50~?}{Oui. $(50 + d)^2 = 2500 + 100d + d^2$, donc on ajoute $d$ à 25 pour les centaines et on écrit $d^2$ sur les deux dernières places. $53^2 = 2809$ et, avec $d = -3$, $47^2 = 2209$.}
\iqfollow{Comment élevez-vous au carré un autre nombre à deux chiffres~?}{On passe au multiple de 10 le plus proche, puis on corrige~: $n^2 = (n + d)(n - d) + d^2$. Donc $67^2 = 70 \times 64 + 9 = 4489$.}
\end{solution}

\begin{solution}{iq:iv1:mental-technique:3}
Le produit vaut 9312. On retranche 3 à 96, ou de façon équivalente 4 à 97, pour obtenir le 93. On multiplie ensuite les écarts à 100, $3 \times 4 = 12$, et on l'écrit sur les deux dernières places.

Soient $a$ et $b$ les écarts~:
\[
(100 - a)(100 - b) = 100\,(100 - a - b) + ab.
\]
Le premier terme est la partie gauche de la réponse, et $ab$ en est la partie droite. La même algèbre vaut pour les nombres supérieurs à 100, où les écarts s'ajoutent~: $103 \times 106 = 100 \times 109 + 18 = 10\,918$.

\iqwrong{Traiter la partie droite comme deux chiffres quoi qu'il arrive. Si $ab$ atteint 100, il y a une retenue. Pour $88 \times 85$, les écarts valent 12 et 15, donc la partie gauche vaut 73 et $ab = 180$, et la réponse est $7300 + 180 = 7480$.}
\iqfollow{Faut-il prendre 100 comme base~?}{Non. Toute base facile convient. Autour de 50, $48 \times 47 = 50 \times 45 + 6 = 2256$.}
\end{solution}

\begin{solution}{iq:iv1:mental-technique:4}
Le produit vaut $46\,000$.

Puisque $125 = 1000/8$, on multiplie par 1000 et on divise par 8. En pratique, on fait d'abord la division, car elle garde les nombres petits~: $368 / 8 = 46$, puis on ajoute trois zéros.

L'astuce consiste à remplacer un multiplicateur difficile par une fraction facile d'une puissance de dix. Les valeurs à retenir sont $5 = 10/2$, $25 = 100/4$, $125 = 1000/8$, $12.5 = 100/8$ et $75 = 300/4$. Si la division ne tombe pas juste, on poursuit en décimales~: $372 / 8 = 46.5$, donc $372 \times 125 = 46\,500$.

\iqfollow{Et pour $368 \times 375$~?}{$375 = 3 \times 125$, donc trois fois la réponse précédente~: $138\,000$.}
\iqfollow{Et pour une division par 125~?}{On fait l'inverse~: on multiplie par 8 et on divise par 1000. $3500 / 125 = 28\,000 / 1000 = 28$.}
\end{solution}

\begin{solution}{iq:iv1:mental-technique:5}
Le produit vaut 5346.

On écrit le premier chiffre, puis la somme de chaque paire de voisins, puis le dernier chiffre, en travaillant depuis la droite pour pouvoir gérer les retenues. Pour 486, les colonnes sont $6$, puis $8 + 6 = 14$, puis $4 + 8 = 12$, puis $4$. En partant de la droite~: on écrit 6~; on écrit 4 et on retient 1~; $12 + 1 = 13$, donc on écrit 3 et on retient 1~; $4 + 1 = 5$. Le résultat est 5346.

La méthode fonctionne parce que $11n = 10n + n$. En écrivant $n$ décalé d'une place au-dessus de $n$ lui-même, chaque chiffre se retrouve dans la même colonne que son voisin, et l'addition en colonnes donne exactement ces sommes~:
\[
4860 + 486 = 5346.
\]

\iqwrong{Oublier la retenue et enchaîner les sommes de colonnes, ce qui donne quelque chose comme 4\,12\,14\,6. Mieux vaut poser l'addition sur deux lignes tant que les retenues ne sont pas automatiques.}
\iqfollow{Quelle est l'astuce équivalente pour 101~?}{On écrit le nombre deux fois, avec un décalage de deux places~: $486 \times 101 = 48\,600 + 486 = 49\,086$.}
\end{solution}

\begin{solution}{iq:iv1:mental-technique:6}
$16~\%$ de 25 vaut 4. Il en va de même pour 25~\% de 16, et c'est justement ce qu'on calcule.

L'échange est légitime, car les deux expressions donnent le même produit divisé par 100~: $p\%$ de $n$ vaut $pn/100$, et la multiplication ne dépend pas de l'ordre. On choisit donc le sens qui rend le pourcentage facile. Ainsi, 18~\% de 50 équivaut à 50~\% de 18, soit 9, et 4~\% de 175 équivaut à 175~\% de 4, soit 7.

\iqwrong{Prendre cela pour une astuce de salon qui ne marcherait qu'avec des nombres ronds. Elle marche toujours. Elle n'aide que lorsque l'un des deux nombres est un pourcentage facile~: 25, 50, 75, 125, ou un multiple de 10. Pour 37~\% de 41, aucun des deux sens n'est facile.}
\iqfollow{Combien valent 36~\% de 25, et 12.5~\% de 64~?}{36~\% de 25 équivaut à 25~\% de 36, soit 9. Et 12.5~\% correspond à un huitième, donc 12.5~\% de 64 vaut 8.}
\end{solution}

\begin{solution}{iq:iv1:mental-technique:7}
La taxe vaut 64, et le prix avant taxe vaut 320.

Le prix toutes taxes comprises représente 120~\% du prix avant taxe, donc la taxe en est $20/120$, soit un sixième, et $384 / 6 = 64$. En général, une taxe de taux $r$ sur le net représente $r/(1 + r)$ du brut.

\iqwrong{Prendre 20~\% de 384, ce qui donne 76.8. Les 20~\% s'appliquent au prix avant taxe, pas au montant qui contient déjà la taxe.}
\iqfollow{Quelle part du brut représente une majoration de 25~\%, puis une de 8~\%~?}{Une majoration de 25~\% représente un cinquième du brut, puisque $25/125 = 1/5$. Une majoration de 8~\% vaut $8/108 = 2/27$, soit environ $7.4~\%$.}
\iqfollow{Vous avez payé 384 après une remise de 20~\%. Quel était le prix catalogue~?}{On a payé 80~\% du prix catalogue, donc il était de $384 / 0.8 = 480$. La remise de 96 représente un quart de ce qu'on a payé, et c'est la même étape inverse.}
\end{solution}

\begin{solution}{iq:iv1:mental-technique:8}
Une hausse de 25~\% est nécessaire. En général, après une baisse de $d$, la hausse nécessaire vaut $d/(1 - d)$. Une baisse de 37.5~\% exige une hausse de 60~\%.

Après la baisse, il reste $1 - d$ du départ, et la hausse doit le ramener à 1, elle doit donc valoir $\frac{1}{1 - d} - 1 = \frac{d}{1 - d}$. Les fractions rendent le calcul rapide~: une baisse d'un cinquième laisse quatre cinquièmes, et il faut gagner un quart. Une baisse de $37.5~\% = 3/8$ laisse $5/8$, et il faut gagner $3/5$. Une baisse de moitié exige une hausse de 100~\%.

\iqwrong{Dire 20~\%. La hausse s'applique à une base plus petite, elle doit donc être plus grande en pourcentage.}
\iqwhy{Une stratégie qui se trouve 20~\% sous son plus haut doit gagner 25~\% pour se rétablir, et une stratégie 50~\% sous doit doubler. Les drawdowns coûtent donc plus qu'il n'y paraît, et ce coût croît vite avec leur taille.}
\end{solution}

\begin{solution}{iq:iv1:mental-technique:9}
$11/14$ est plus grand, de $1/126$.

Par multiplication en croix\index{multiplication en croix}~: $a/b$ et $c/d$ se comparent comme $ad$ et $bc$, pourvu que les dénominateurs soient positifs. Ici $7 \times 14 = 98$ et $9 \times 11 = 99$, donc la seconde fraction est plus grande. L'écart est la différence des produits croisés divisée par le produit des dénominateurs~:
\[
\frac{11}{14} - \frac{7}{9} = \frac{99 - 98}{14 \times 9} = \frac{1}{126}.
\]
Une seconde méthode compare les distances à 1~: $2/9 \approx 0.222$ et $3/14 \approx 0.214$, et le plus petit écart revient à la plus grande fraction.

\iqwrong{Dire que la fraction aux plus grands nombres doit être la plus grande. Cela ne prouve rien~: $1/2$ dépasse $11/30$.}
\iqfollow{Lequel est le plus grand, $13/17$ ou $23/30$~?}{$13 \times 30 = 390$ et $17 \times 23 = 391$, donc $23/30$ est plus grand, de $1/510$.}
\end{solution}

\begin{solution}{iq:iv1:mental-technique:10}
$3/7 = 0.428571\dots$ et $5/7 = 0.714285\dots$. On garde les mêmes six chiffres dans le même ordre cyclique, et on part du chiffre qui donne la bonne grandeur.

Pour $3/7$, on sait qu'il vaut environ $0.43$, donc le cycle doit commencer par 4~: 428571. Pour $5/7$, environ $0.71$, il commence par 7~: 714285.

La raison tient à la division longue. Diviser 1 par 7 passe par les restes $1, 3, 2, 6, 4, 5$ puis revient à 1, soit tous les restes non nuls possibles. Partir de 3 au lieu de 1 revient seulement à rejoindre cette même boucle en un autre point, donc on lit les mêmes chiffres à partir d'un autre endroit.

\iqfollow{Toute fraction se comporte-t-elle de cette façon~?}{Non. Il faut que les restes forment une seule boucle. Pour 13, les restes se répartissent en deux boucles, donc $1/13 = 0.076923\dots$ et $3/13 = 0.230769\dots$ partagent un même cycle, mais $2/13 = 0.153846\dots$ n'en fait pas partie.}
\end{solution}

\begin{solution}{iq:iv1:mental-technique:11}
Le prix moyen est de 85.

Utiliser les écarts à une base\index{moyenne pondérée par les écarts}. Prendre 80 comme base. Les 500 actions à 88 sont à 8 au-dessus, et elles représentent $5/8$ du total, donc la moyenne se situe aux $5/8$ du chemin de 80 à 88. Cela donne $80 + 5 = 85$.

Une moyenne pondérée de deux nombres se situe toujours entre eux, à une distance proportionnelle au poids du côté opposé. Il suffit donc du rapport des poids, jamais des totaux.

\iqfollow{Vous détenez 400 à $101.5$, 100 à 102 et 500 à 101. Quel est votre prix moyen~?}{On prend 101 comme base. Les écarts valent $0.5$, 1 et 0, avec des poids $0.4$, $0.1$ et $0.5$. Leur somme pondérée vaut $0.2 + 0.1 = 0.3$, donc le prix moyen est $101.3$.}
\iqwrong{Faire la moyenne des cours sans pondération~: $(80 + 88)/2 = 84$. Le plus gros ordre attire la moyenne vers son propre cours, donc la réponse doit être supérieure à 84.}
\end{solution}

\begin{solution}{iq:iv1:mental-technique:12}
La règle donne 9 ans, et la valeur exacte est $9.006$ ans. L'approximation est si précise à 8~\%.

Résoudre $(1 + r)^t = 2$, ce qui donne $t = \ln 2 / \ln(1 + r)$. Pour $r$ petit, $\ln(1 + r) \approx r$, donc $t \approx 0.693/r$, soit $69.3$ divisé par le taux exprimé en pourcentage. Or $\ln(1 + r)$ est un peu inférieur à $r$, ce qui pousse le numérateur à la hausse~: à 8~\%, le numérateur exact vaut environ $72.05$. Le nombre 72 est aussi choisi parce qu'il est divisible par 2, 3, 4, 6, 8, 9 et 12.

\iqwrong{Raisonner avec l'intérêt simple~: $100/8 = 12.5$ ans. Cela ignore les intérêts sur les intérêts, ce qui fait tout l'enjeu des règles de doublement.}
\iqfollow{Combien de temps faut-il pour tripler à 8~\%~?}{Utilisons $\ln 3 \approx 1.1$ de la même façon. À ce taux, le numérateur monte à environ 114, donc la réponse est d'environ $114/8 = 14.25$ ans. La valeur exacte vaut $14.27$.}
\end{solution}

\begin{solution}{iq:iv1:mental-technique:13}
Environ 2~\% par jour. Le chiffre rapide est 16.

Les variances des rendements indépendants s'additionnent, donc l'écart type croît comme la racine carrée du temps\index{règle de la racine carrée du temps}. La volatilité annuelle vaut donc la volatilité quotidienne multipliée par $\sqrt{252} \approx 15.87$. Utiliser 16 ne commet qu'une erreur d'environ 1~\%, ce qui est bien meilleur que n'importe quelle estimation de la volatilité elle-même. Ainsi $32\% / 16 = 2\%$~; la valeur exacte est $2.02\%$.

Les deux autres nombres utiles à retenir sont $\sqrt{52} \approx 7.2$ pour les semaines et $\sqrt{12} \approx 3.46$ pour les mois. Avec une volatilité annuelle de 32~\%, l'écart type du rendement d'une semaine vaut environ $4.4~\%$, et celui d'un mois environ $9.2~\%$.

\iqwrong{Diviser par 252, ce qui donne $0.13~\%$. Cela traite la volatilité comme si elle s'additionnait dans le temps, comme le fait un rendement. Ce sont les variances qui s'additionnent, et non les volatilités.}
\iqwhy{C'est le réflexe derrière chaque chiffre de risque quotidien. En voyant une action gagner 4~\% en une journée avec une volatilité annuelle de 32~\%, on sait tout de suite qu'il s'agit d'une journée à deux écarts types.}
\end{solution}

\begin{solution}{iq:iv1:mental-technique:14}
$\log_{10} 72 \approx 1.857$ et $\log_{10} 0.0072 \approx -2.143$.

Décomposer le nombre en facteurs premiers connus, et transformer le produit en somme\index{logarithmes par factorisation}. Puisque $72 = 2^3 \times 3^2$,
\[
\log_{10} 72 = 3(0.301) + 2(0.477) = 0.903 + 0.954 = 1.857.
\]
La valeur exacte est $1.8573$. Pour le second, $0.0072 = 72 \times 10^{-4}$, donc on retranche 4 et on obtient $-2.143$.

On dispose aussi de $\log_{10} 5$ sans calcul supplémentaire~: $5 = 10/2$, donc il vaut $1 - 0.301 = 0.699$. Ces trois valeurs, avec les puissances de 10, couvrent tout nombre formé des nombres premiers 2, 3 et 5. C'est pour cela que ce couple mérite d'être mémorisé.

\iqfollow{Que vaut $\log_{10} 1.5$~?}{$1.5 = 3/2$, donc sa valeur est $0.477 - 0.301 = 0.176$.}
\iqwrong{Prendre le log de $0.0072$ comme un nombre positif compris entre 0 et 1. Un nombre inférieur à 1 a un log négatif, donc un contrôle du signe repère cette erreur tout de suite.}
\end{solution}

\begin{solution}{iq:iv1:mental-technique:15}
$2^{30}$ vaut environ $10^9$, soit un milliard. Mille milliards demandent environ 40 doublements. La première estimation est d'environ 7~\% trop basse.

Le fait utile est que $2^{10} = 1024 \approx 10^3$. Donc $2^{30} = (2^{10})^3 \approx 10^9$, et $2^{40} \approx 10^{12}$. La vraie valeur de $2^{30}$ est $1\,073\,741\,824$, soit $1.024^3 \approx 1.074$ milliard. L'erreur s'accumule à raison d'environ $2.4~\%$ par facteur de 1024, donc elle pèse davantage pour les grandes puissances~: $2^{100}$ vaut environ $1.27 \times 10^{30}$ et non $10^{30}$.

\iqfollow{Quel est l'ordre de grandeur de $2^{64}$~?}{$2^{60} \approx 1.15 \times 10^{18}$, en utilisant $1.024^6 \approx 1.15$, et $2^4 = 16$, donc environ $1.8 \times 10^{19}$.}
\iqwhy{Les développeurs lisent ainsi les ordres de grandeur en permanence. Une table d'un milliard de clés demande environ 30 bits par clé pour les distinguer, et elle dépasse vite la plage des 32 bits.}
\end{solution}

\begin{solution}{iq:iv1:mental-technique:16}
$\sqrt{50} \approx 7.071$. L'estimation $7 + 1/14 = 7.0714$ est légèrement trop élevée, et la valeur exacte est $7.07107$.

Prendre le carré facile le plus proche, $a^2 = 49$, avec $d = 50 - 49 = 1$, et utiliser
\[
\sqrt{a^2 + d} \approx a + \frac{d}{2a}.
\]
Il s'agit d'une étape de la méthode de Newton\index{méthode de Newton}. Élever l'estimation au carré montre pourquoi elle dépasse toujours la valeur visée~: $(a + d/2a)^2 = a^2 + d + (d/2a)^2$, qui est au-dessus de la cible d'une quantité positive, quel que soit le signe de $d$. Une correction supplémentaire, toujours vers le bas, la rapproche beaucoup, et elle reste juste au-dessus de la valeur vraie.

\iqfollow{Estimez $\sqrt{3\,200\,000}$.}{On l'écrit sous la forme $1000 \sqrt{3.2}$. Comme $1.8^2 = 3.24$, on prend $a = 1.8$ et $d = -0.04$. Cela donne $1.8 - 0.04/3.6 = 1.789$, soit environ 1789. La valeur exacte est $1788.9$.}
\iqwrong{Interpoler linéairement entre $\sqrt{49}$ et $\sqrt{64}$, ce qui donne $7 + 1/15$. La pente est trop plate pour cette partie de la courbe, et le résultat sort trop bas d'environ $0.004$.}
\end{solution}

\begin{solution}{iq:iv1:mental-technique:17}
10 points de base sur 3.5 millions de dollars font 3500~\$, et l'intérêt supplémentaire sur trois mois vaut 6250~\$.

Retenir un point de repère~: un point de base sur 1 million de dollars vaut 100~\$. Donc $n$ points de base sur $m$ millions de dollars font $100\,nm$ dollars~: $100 \times 10 \times 3.5 = 3500$. Pour le prêt, un point de base sur 250 millions de dollars vaut 25\,000~\$ par an, et trois mois en représentent un quart, soit 6250~\$.

\iqwrong{Confondre un point de base avec $0.01$, qui vaut un pour cent entier. En décimal, un point de base vaut $0.0001$. Si la réponse pour le prêt ressemble à 625\,000~\$, on a commis cette erreur.}
\iqfollow{Combien de points de base représentent un huitième de pour cent~?}{12.5, puisque 1~\% vaut 100 points de base.}
\iqwhy{Les intérêts, les spreads et les frais se cotent tous en points de base, et la conversion en dollars doit venir sans réflexion.}
\end{solution}

\begin{solution}{iq:iv1:mental-technique:18}
99-17 vaut $99.53125$.

Retenir que $1/32 = 0.03125$, en divisant successivement par deux à partir de un demi~: $0.5$, $0.25$, $0.125$, $0.0625$, $0.03125$. Puis on construit la fraction à partir de morceaux~: $17/32 = 16/32 + 1/32 = 0.5 + 0.03125 = 0.53125$.

\iqfollow{Dans l'autre sens, comment coteriez-vous $101.40625$~?}{La partie fractionnaire vaut $0.40625$. On retire $0.375$, soit $12/32$, et il reste $0.03125$, soit $1/32$. On obtient donc $13/32$, soit 101-13.}
\iqfollow{Quelle est la valeur en dollars d'un mouvement de prix d'un 32\up{e} sur 10 millions de dollars de valeur nominale~?}{Le prix est exprimé pour 100 de valeur nominale, donc un 32\up{e} vaut $0.03125$ pour 100, soit $0.03125~\%$ de la valeur nominale, soit 3125~\$.}
\iqwrong{Lire 99-17 comme $99.17$. Les chiffres après le tiret comptent des trente-deuxièmes, et non des centièmes~: le prix est donc de plus d'un demi-point au-dessus de 99, et non à environ un sixième.}
\end{solution}

\begin{solution}{iq:iv1:mental-technique:19}
$38 \times 42$ est plus grand, exactement de 5.

Les deux couples ont pour somme 80, donc pour moyenne 40. On écrit chacun sous la forme $40$ plus ou moins son demi-écart~:
\[
37 \times 43 = 40^2 - 3^2 = 1591, \qquad 38 \times 42 = 40^2 - 2^2 = 1596.
\]
La règle générale~: à somme fixée, le produit est le carré de la moyenne moins le carré du demi-écart\index{difference@différence de deux carrés}. Ainsi, plus les nombres sont proches, plus leur produit est grand. Il est maximal lorsqu'ils sont égaux, et chaque pas d'écart coûte le carré du demi-écart, et non une quantité fixe.

\iqfollow{Deux nombres ont pour somme 100. Quel est le plus grand produit possible~?}{Cinquante chacun donne 2500. Des écarts de 2, 10 et 20 donnent 2499, 2475 et 2400, donc la perte est faible près du sommet et croît vite à mesure qu'on s'en éloigne.}
\iqfollow{La règle s'étend-elle à davantage de facteurs~?}{Oui. À somme fixée, le produit de nombres positifs est maximal lorsqu'ils sont égaux. C'est l'inégalité des moyennes arithmétique et géométrique\index{inégalité des moyennes arithmétique et géométrique}.}
\end{solution}

\begin{solution}{iq:iv1:mental-technique:20}
$1.02 \times 0.97 = 0.9894$. Pour une réponse approximative, on néglige le terme croisé, ce qui laisse $0.99$, avec une erreur de $0.0006$.

On développe avec $a = 0.02$ et $b = -0.03$~:
\[
(1 + a)(1 + b) = 1 + a + b + ab = 1 + 0.02 - 0.03 - 0.0006 = 0.9894.
\]
La partie du premier ordre, $1 + a + b$, est la somme des deux mouvements. Le terme croisé $ab$ est le produit de deux petits nombres, il est donc beaucoup plus petit, et son signe est le produit de leurs signes.

\iqwrong{Se tromper sur le signe du terme croisé et répondre $0.9906$. Avec un mouvement à la hausse et un à la baisse, le terme croisé est négatif.}
\iqfollow{Et $1.02 \times 1.03 \times 0.96$~?}{Au premier ordre~: $1 + 0.02 + 0.03 - 0.04 = 1.01$. Les produits deux à deux valent $0.0006 - 0.0008 - 0.0012 = -0.0014$, et le produit triple vaut environ $-0.00002$. On obtient $1.00858$, et la valeur exacte est $1.008576$.}
\iqwhy{Les rendements sur quelques périodes s'additionnent au premier ordre. On sait ainsi mesurer l'ampleur d'une telle erreur.}
\end{solution}

\begin{solution}{iq:iv1:mental-technique:21}
La réponse est fausse, car elle échoue au contrôle par la somme des chiffres. Elle passe en revanche les contrôles du dernier chiffre et de l'ordre de grandeur. Le produit exact vaut $100\,022$.

Dernier chiffre~: $7 \times 6 = 42$ se termine par 2, comme la réponse proposée. Ordre de grandeur~: $3850 \times 26 = 100\,100$ est très proche de la réponse proposée. Aucun des deux contrôles ne permet donc de trancher.

Preuve par neuf\index{preuve par neuf}~: un nombre a le même reste modulo 9 que la somme de ses chiffres, car $10 \equiv 1$ modulo 9. Ici, 3847 a pour somme des chiffres 22, qui se réduit à 4. Et 26 se réduit à 8. Le produit doit donc se réduire à $4 \times 8 = 32$, qui se réduit à 5. La somme des chiffres de la réponse proposée vaut 6, et non 5~: elle est donc fausse.

\iqfollow{Si le contrôle passe, le produit est-il juste~?}{Non. Un succès indique seulement que l'erreur est un multiple de 9, ce qui, pour une erreur au hasard, arrive environ une fois sur neuf. Le contrôle peut prouver qu'un produit est faux, mais pas qu'il est juste.}
\iqwhy{Cela ne prend que quelques secondes et détecte la plupart des étourderies, y compris un chiffre faux, qui est l'erreur habituelle. On l'utilise sur ses propres calculs avant de les énoncer.}
\end{solution}

\begin{solution}{iq:iv1:mental-technique:22}
C'est 4964, retenu par la preuve par neuf\index{preuve par neuf}.

Chaque option se termine par 4, comme il se doit, puisque $3 \times 8 = 24$. Le dernier chiffre ne permet donc pas de trancher. L'ordre de grandeur ne tranche pas non plus facilement, car $70 \times 70 = 4900$ indique seulement que la réponse se situe dans la bonne région. Il faut donc comparer les restes modulo 9. 73 a pour somme des chiffres 10, qui se réduit à 1, et 68 a pour somme 14, qui se réduit à 5. Le produit doit donc se réduire à 5. Les options se réduisent à $4964 \to 23 \to 5$, $4954 \to 22 \to 4$, $5064 \to 15 \to 6$ et $4864 \to 22 \to 4$. Seule la première convient.

Pour le vérifier directement, la différence de deux carrés fonctionne avec des demi-entiers~: la moyenne vaut $70.5$ et l'écart à la moitié vaut $2.5$, donc $70.5^2 - 2.5^2 = 4970.25 - 6.25 = 4964$.

\iqwrong{Choisir l'option la plus proche de $70 \times 70 = 4900$. Cela donne 4864, qui est faux. Un produit arrondi indique où se situe la réponse, et non laquelle de plusieurs options voisines est la bonne.}
\end{solution}

\begin{solution}{iq:iv1:mental-technique:23}
Un contrôle par la somme des chiffres ne détecte pas cette erreur, car 6428 est 6248 avec deux chiffres échangés, et échanger des chiffres laisse la somme inchangée. On peut utiliser un contrôle modulo 11, ou le fait que 8 divise 88. Le produit exact vaut 6248.

Les sommes des chiffres concordent~: 71 se réduit à 8, 88 se réduit à 7, et $8 \times 7 = 56$ se réduit à 2~; et 6428 a pour somme des chiffres 20, qui se réduit à 2. Le contrôle par neuf passe donc.

Avec onze, $10 \equiv -1$ modulo 11\index{preuve par onze}, donc un nombre a le même reste que la somme alternée de ses chiffres, lue depuis la droite. Comme $88 = 8 \times 11$, le produit doit laisser le reste 0. Or $8 - 2 + 4 - 6 = 4$, donc 6428 échoue. Le vrai 6248 a $8 - 4 + 2 - 6 = 0$. Le facteur 8 la détecte aussi, car $428 / 8$ n'est pas entier, alors que $248 / 8 = 31$.

\iqfollow{Pourquoi le contrôle par onze détecte-t-il les échanges~?}{Échanger deux chiffres voisins place chacun dans une position de signe opposé, de sorte que la somme alternée change de deux fois leur différence, et cela n'est pas un multiple de 11 sauf si les chiffres sont égaux.}
\end{solution}

\begin{solution}{iq:iv1:mental-technique:24}
Le changement net vaut $+2~\%$, et l'ordre n'a pas d'importance.

On multiplie les facteurs~: $1.2 \times 0.85 = 1.02$. La multiplication ne dépend pas de l'ordre. La règle pour composer deux mouvements est $a + b + ab$~:
\[
0.20 - 0.15 + (0.20)(-0.15) = 0.05 - 0.03 = 0.02.
\]
On additionne donc les deux mouvements, puis on ajoute leur produit en surveillant son signe. Ce produit porte toujours sur les deux mouvements pris ensemble, il est donc négatif lorsqu'ils sont de signes opposés.

\iqwrong{Répondre $+5~\%$, en additionnant les deux pourcentages. À partir de 100, la hausse ajoute 20, et la baisse, calculée sur la base plus élevée de 120, retire 18 au lieu de 15.}
\iqfollow{Et une hausse de 10~\% suivie d'une baisse de 10~\%~?}{$-1~\%$~: les mouvements se compensent et seul le terme croisé subsiste. De même, $+50~\%$ suivi de $-50~\%$ laisse un recul de $25~\%$.}
\end{solution}

\begin{solution}{iq:iv1:mental-technique:25}
$1/0.97 \approx 1.0309$ et $1/1.03 \approx 0.9709$.

On utilise la série géométrique~: pour $x$ petit,
\[
\frac{1}{1 - x} = 1 + x + x^2 + \dots, \qquad \frac{1}{1 + x} = 1 - x + x^2 - \dots.
\]
Avec $x = 0.03$, les termes valent $0.03$, $0.0009$ et $0.000027$, donc $1/0.97 \approx 1.030927$ et $1/1.03 \approx 0.970873$. Le terme suivant est inférieur à $10^{-6}$, donc trois termes sont exacts à environ six décimales près.

La méthode cesse de fonctionner quand $x$ augmente~: chaque terme n'est que $x$ fois le précédent, donc pour $x = 0.2$ il faut beaucoup de termes. Dans ce cas, on passe aux fractions~: $1/0.8 = 5/4 = 1.25$.

\iqwrong{Répondre $1/0.97 \approx 1.03$ en s'arrêtant là. C'est juste à deux décimales, mais la question en demande quatre, et le terme carré les fournit.}
\iqwhy{C'est la méthode pour inverser à la volée un taux de change ou un facteur d'actualisation. Un taux de $0.97$ dans un sens vaut $1.0309$ dans l'autre, et non $1.03$. Sur un gros ordre, ces $0.0009$ manquants correspondent à 9 points de base.}
\end{solution}

\begin{solution}{iq:iv1:mental-technique:26}
$1.05^{10} \approx 1.63$. La valeur exacte est $1.62889$.

Première méthode, par élévations au carré successives~: $1.05^2 = 1.1025$. L'élévation au carré donne $1.05^4 \approx 1.2155$. On multiplie par $1.05$ pour obtenir $1.05^5 \approx 1.2763$. En l'élevant au carré, on obtient $1.629$. Les carrés paraissent difficiles mais sont faciles~: $1.1025^2 \approx 1 + 2(0.1025) + 0.1025^2 \approx 1.2155$.

Deuxième méthode, par les logarithmes\index{approximation exponentielle}~: $\ln 1.05 \approx 0.05 - 0.00125 + 0.00004 = 0.04879$, donc dix fois ce nombre donne $0.4879$. On a alors $e^{0.4879} = e^{0.5} e^{-0.0121} \approx 1.6487 \times 0.988 \approx 1.629$.

\iqwrong{Répondre $e^{0.5} \approx 1.65$ en utilisant $\ln(1 + x) \approx x$. Comme $\ln(1 + x)$ est inférieur à $x$, ce calcul donne toujours une valeur trop élevée. L'autre mauvaise réponse est le simple $1.50$, qui ignore les intérêts composés.}
\iqfollow{Que dit la règle de 72~?}{Le doublement prend $72/5 = 14.4$ ans, donc dix ans donnent $2^{10/14.4} = 2^{0.694} \approx 1.62$. C'est plus approximatif, mais cela ne prend que quelques secondes.}
\iqfollow{Comment pouvez-vous le vérifier~?}{On utilise le développement binomial~: $1 + 10(0.05) + 45(0.0025) + 120(0.000125) + 210(0.00000625) \approx 1.6288$. Les termes décroissent rapidement, donc cinq suffisent.}
\end{solution}

\begin{solution}{iq:iv1:mental-technique:27}
$0.99^{100} \approx 0.366$. La valeur exacte est $0.36603$.

On prend les logarithmes\index{approximation exponentielle}~: $\ln 0.99 = -0.01 - 0.00005 - \dots \approx -0.01005$. Donc $100 \ln 0.99 \approx -1.005$, et
\[
0.99^{100} \approx e^{-1} e^{-0.005} \approx 0.3679 \times 0.995 \approx 0.366.
\]

La réponse $1/e \approx 0.368$ utilise $\ln(1 - x) \approx -x$ et néglige le terme suivant de la série, $-x^2/2$. Sur 100 périodes, le terme manquant vaut $-100 \times 0.00005 = -0.005$, soit une variation d'environ un demi pour cent qui se voit à la troisième décimale.

\iqfollow{Et combien vaut $1.01^{100}$~?}{De la même façon, $100 \ln 1.01 \approx 0.995$, donc $e \times 0.995 \approx 2.705$. La valeur exacte est $2.70481$. Notons que $0.99^{100}$ et $1/1.01^{100} = 0.36971$ sont des nombres différents, et que tous deux diffèrent de $1/e$.}
\iqwhy{C'est la probabilité qu'aucun événement ne survienne en 100 essais indépendants à 1~\% chacun. Ainsi, 100 essais à 1~\% produisent au moins un événement avec une probabilité d'environ $63.4~\%$, et non 100~\%.}
\end{solution}

\begin{solution}{iq:iv1:mental-technique:28}
Pour $r = 1$, la règle donne 72 ans contre une valeur exacte de $69.66$, soit un écart de $3.4~\%$ par excès. Pour $r = 36$, elle donne 2 ans contre une valeur exacte de $2.25$, soit $11~\%$ par défaut. Elle est exacte autour de $r = 7.85$. Une bonne correction consiste à prendre un numérateur égal à $69.3 + r/3$.

Le temps de doublement exact vaut $\ln 2 / \ln(1 + r/100)$. En l'écrivant sous la forme d'un numérateur sur $r$, le numérateur vaut $r \ln 2 / \ln(1 + r/100)$, qui n'est pas constant~: il croît avec $r$ à partir de $69.3$. Ses valeurs valent environ $69.66$ pour $r = 1$, $72.05$ pour $r = 8$ et $81.2$ pour $r = 36$.

En développant $\ln(1 + x)$ au deuxième ordre, on obtient $69.3\,(1 + x/2)$ avec $x = r/100$, soit une pente de $0.35$ par point de pourcentage. Une pente de $1/3$ est plus commode à utiliser. Pour $r = 36$, cela donne $81.3/36 = 2.26$, et pour $r = 1$ on obtient $69.6$.

\iqwrong{Croire que 72 est la constante théorique. La constante théorique vaut $69.3$. Le nombre 72 est un nombre commode qui s'accorde par hasard aux taux usuels.}
\end{solution}

\begin{solution}{iq:iv1:mental-technique:29}
Un rendement total d'environ $28.6~\%$. La réponse simple, $25.2~\%$, est trop faible car elle ignore les intérêts composés.

On convertit en taux continu\index{approximation exponentielle}. $\ln 1.001 \approx 0.001 - 0.0000005 = 0.0009995$, et $252$ de ces valeurs donnent $0.2519$. On obtient alors
\[
e^{0.2519} \approx 1 + 0.2519 + \tfrac12 (0.2519)^2 + \tfrac16 (0.2519)^3 \approx 1.2863,
\]
et le terme suivant de la série ajoute $0.00017$. La valeur exacte de $1.001^{252}$ vaut $1.28643$, donc le rendement vaut $28.6~\%$.

\iqfollow{Et si vous perdez 10 points de base par jour à la place~?}{$0.999^{252} \approx e^{-0.2521} \approx 0.777$, soit une perte de $22.3~\%$. Le gain et la perte ne sont donc pas de même ampleur, et les deux ensemble donnent $1.2864 \times 0.7771 < 1$.}
\iqfollow{Et sur cinq ans, au même taux quotidien~?}{Le taux dans l'exposant est cinq fois plus grand, environ $1.26$, donc le facteur vaut $e^{1.26} \approx 3.5$. Ici, la somme simple de $126~\%$ est bien trop faible.}
\end{solution}

\begin{solution}{iq:iv1:mental-technique:30}
Il reste environ $60.5~\%$ de votre argent, soit une perte de $39.5~\%$.

Une hausse suivie d'une baisse multiplie votre argent par $1.1 \times 0.9 = 0.99$, et il y a 50 paires de ce type, donc le résultat vaut $0.99^{50}$. En prenant les logarithmes, $50 \ln 0.99 \approx -0.5025$, et $e^{-0.5025} \approx 0.605$.

La moyenne arithmétique des rendements est nulle, mais l'argent se compose, et il se compose au taux géométrique\index{traînée de volatilité}. La règle est que le taux de croissance par période vaut à peu près la moyenne moins la moitié de la variance, $\mu - \sigma^2/2$. Ici $\sigma = 0.1$, donc la pénalité vaut $0.005$ par période, soit $0.5$ sur 100 périodes, et $e^{-0.5} \approx 0.607$.

\iqwrong{Dire que vous revenez au point de départ parce que les variations se compensent. Elles se compensent en addition, mais les rendements se multiplient.}
\iqfollow{Et avec des hausses et des baisses de 20~\%~?}{$0.96^{50} \approx 0.130$, à comparer avec $e^{-2} \approx 0.135$ donné par la règle. La règle devient approximative, car les termes d'ordre supérieur ont maintenant de l'importance.}
\iqwhy{Une volatilité quotidienne de 1~\% coûte environ $0.01^2/2 \times 252 = 1.26~\%$ par an de croissance composée. Doubler la volatilité et le coût est quatre fois plus grand, car il varie avec le carré.}
\end{solution}

\begin{solution}{iq:iv1:mental-technique:31}
Environ 23 ans, plus précisément $22.76$. La règle de 72 donne 24, ce qui est trop long.

Le temps exact résout $0.97^t = 0.5$, donc $t = \ln 2 / (-\ln 0.97)$. En développant, $-\ln(1 - x) = x + x^2/2 + \dots$ est un peu supérieur à $x$, donc le dénominateur vaut environ $0.0305$ et $t$ est inférieur à $69.3/3 = 23.1$.

Le doublement se comporte à l'inverse. Là, $\ln(1 + x) = x - x^2/2 + \dots$ est un peu inférieur à $x$, donc le temps de doublement à 3~\% est plus long, soit $23.45$ ans. À taux égal, il faut donc plus de temps pour doubler qu'il n'en faut pour diviser par deux. La règle de 72 est trop longue dans les deux cas à 3~\%.

\iqfollow{Existe-t-il un numérateur corrigé pour la décroissance~?}{Oui. Là où le doublement utilise $69.3 + r/3$, la décroissance utilise $69.3 - r/3$. À 3~\%, cela donne $68.3/3 = 22.8$ ans.}
\iqwrong{Diviser 50~\% par 3~\%, ce qui donne $16.7$ ans. Cela applique 3~\% du montant initial chaque année, alors qu'ici les 3~\% s'appliquent à ce qui reste, et les pertes diminuent donc.}
\end{solution}

\begin{solution}{iq:iv1:mental-technique:32}
$\log_{10} 7 \approx 0.845$. La valeur exacte est $0.8451$.

Chercher une puissance de 7 proche d'un nombre formé uniquement de 2, 3 et 5\index{logarithmes par factorisation}. Comme $7^4 = 2401$, et $2400 = 2^5 \times 3 \times 5^2$,
\[
\log_{10} 7 \approx \tfrac14 \log_{10} 2400 = \tfrac14 \bigl[ 5(0.301) + 0.477 + 2(0.699) \bigr] = \tfrac14 (3.380) = 0.845,
\]
en utilisant $\log_{10} 5 = 1 - 0.301 = 0.699$. Le rapport $2401/2400$ ne diffère de 1 que de $1/2400$, et diviser par quatre rend l'erreur encore plus petite, donc l'arrondi des données est la seule véritable erreur.

Une voie plus grossière utilise $7^2 = 49 \approx 50$, ce qui donne $(1.699)/2 = 0.8495$, faux à la troisième décimale.

\iqwrong{Interpoler linéairement entre $\log_{10} 6 = 0.778$ et $\log_{10} 8 = 0.903$, ce qui donne $0.8405$. La courbe du logarithme est incurvée, donc une droite passant par deux points de cette courbe se situe sous la courbe.}
\iqfollow{Et $\log_{10} 13$~?}{Utiliser la même idée avec $13 \times 5 = 65$, voisin de $64 = 2^6$. Ainsi $\log_{10} 13 = \log_{10} 64 + \log_{10}(65/64) - \log_{10} 5$. Avec $\log_{10}(65/64) \approx 0.0155 \times 0.434 \approx 0.0067$, cela vaut $1.806 + 0.0067 - 0.699 \approx 1.114$.}
\end{solution}

\begin{solution}{iq:iv1:mental-technique:33}
$2^{100}$ a 31 chiffres, et son premier chiffre est 1. Il vaut environ $1.27 \times 10^{30}$.

Prendre les logarithmes~: $100 \times 0.30103 = 30.103$. La partie entière indique que le nombre est compris entre $10^{30}$ et $10^{31}$, donc il a 31 chiffres. La partie fractionnaire donne les premiers chiffres~: $10^{0.103}$. Comme $\log_{10} 1.25 = 1 - 3(0.30103) = 0.0969$, le reste $0.103 - 0.0969 = 0.006$ se multiplie par environ $1.014$, ce qui donne $1.25 \times 1.014 \approx 1.268$. Le nombre exact commence par $1.2676$.

\iqfollow{Et pour $3^{50}$~?}{$50 \times 0.47712 = 23.856$, donc 24 chiffres. Comme $\log_{10} 7 \approx 0.845$, la partie principale est un peu supérieure à 7. Elle vaut $7.18 \times 10^{23}$.}
\iqwrong{Compter les chiffres comme la partie entière du logarithme, ici 30. Un nombre de logarithme $30.103$ a 31 chiffres, tout comme $10^{30}$ lui-même en a 31.}
\iqwhy{Savoir évaluer la taille d'une puissance permet de voir immédiatement si un compte tient dans un entier de 64 bits. Comme $2^{64}$ a 20 chiffres, un compte de 21 chiffres ne tient pas.}
\end{solution}

\begin{solution}{iq:iv1:mental-technique:34}
$\ln 1.3 \approx 0.262$. La valeur exacte est $0.26236$.

Trouvez un nombre voisin dont vous connaissez le logarithme, puis faites un petit pas à partir de lui. Ici $1.25 = 10/8$, donc
\[
\ln 1.25 = \ln 10 - 3 \ln 2 = 2.303 - 2.079 = 0.224.
\]
Ensuite $1.3 / 1.25 = 1.04$, et $\ln 1.04 \approx 0.04 - 0.0008 = 0.0392$. Le total vaut $0.263$, et en utilisant davantage de chiffres pour les deux logarithmes donnés, on obtient $0.262$.

La série $\ln(1 + x) = x - x^2/2 + x^3/3 - \dots$ avec $x = 0.3$ converge lentement~: les trois premiers termes valent $0.3$, $-0.045$ et $+0.009$. En s'arrêtant après deux termes on obtient $0.255$, et après trois $0.264$. Un petit pas depuis un point voisin connu ne demande qu'un ou deux termes.

\iqfollow{Retrouvez $\ln 3$ à partir de ces deux mêmes logarithmes.}{Comme $3^4 = 81$ est très proche de $80 = 2^3 \times 10$, on obtient $4 \ln 3 \approx 3(0.693) + 2.303 + \ln(81/80) = 4.382 + 0.0124$, donc $\ln 3 \approx 1.099$.}
\iqwhy{Un rendement simple de 30~\% correspond à un rendement logarithmique d'environ $26.2~\%$. L'écart explique pourquoi des rendements de cette taille ne peuvent pas simplement s'additionner.}
\end{solution}

\begin{solution}{iq:iv1:mental-technique:35}
La racine dixième de 2 vaut environ $1.072$, et la racine cinquième de 3 environ $1.246$.

Une racine devient une division dans les logarithmes\index{approximation exponentielle}. Pour la première, $\ln 2 / 10 = 0.0693$, puis
\[
e^{0.0693} \approx 1 + 0.0693 + \tfrac12 (0.0693)^2 + \tfrac16 (0.0693)^3 \approx 1.0718.
\]
Pour la seconde, $\ln 3 / 5 = 0.2198$, et $1 + 0.2198 + 0.0242 + 0.0018 + 0.0001 \approx 1.246$. La valeur exacte est $1.2457$.

Cela se relie à la règle de 72~: la racine dixième de 2 signifie que doubler en dix ans demande $7.18~\%$ par an, et la règle donne $7.2~\%$.

\iqwrong{Diviser l'excédent par 10 et dire $1.1$. C'est un intérêt simple, et non composé. Un taux de croissance de 10~\% par an pendant dix ans donne $2.59$, et non 2.}
\iqfollow{Quel taux mensuel se compose en 12~\% par an~?}{$\ln 1.12 \approx 0.1133$, et diviser par 12 donne $0.00944$, soit environ $0.949~\%$ par mois, et non 1~\%.}
\end{solution}

\begin{solution}{iq:iv1:mental-technique:36}
Le produit vaut environ $4.47$ millions, à moins de $0.3~\%$ de la valeur exacte de $4\,484\,646$.

Arrondir à des nombres faciles à multiplier~: $50 \times 60 \times 80 \times 20 = 4\,800\,000$. Corriger ensuite à l'aide de chaque arrondi exprimé en pourcentage. Les valeurs arrondies sont trop fortes de $6.4~\%$ ($50/47$) et de $5.3~\%$ ($20/19$), et trop faibles de $3.2~\%$ ($60/62$) et de $1.2~\%$ ($80/81$). Pour de petites erreurs, elles s'additionnent, donc le total vaut environ $+7.3~\%$. Diviser $4.8$ millions par $1.073$ pour obtenir environ $4.47$ millions.

Un bon arrondi associe un nombre arrondi par excès à un autre arrondi par défaut, de sorte que les corrections se compensent en partie et restent assez petites pour s'additionner. Vous avez procédé ainsi sans une seule multiplication longue.

\iqwrong{Arrondir les quatre facteurs et s'arrêter à $4.8$ millions. Cela fait 7~\% d'écart, bien que le résultat semble correct parce que chaque facteur paraît proche. Quatre petites erreurs s'additionnent, il faut donc les corriger.}
\iqfollow{Et si les erreurs ne sont pas petites~?}{Multiplier les rapports au lieu de les additionner, ou additionner les logarithmes. Additionner des pourcentages ne convient que tant que chacun reste dans un ordre de grandeur d'environ 10~\%.}
\end{solution}

\begin{solution}{iq:iv1:mental-technique:37}
La probabilité d'obtenir exactement 10 piles est d'environ $0.18$ (exactement $0.1762$), et celle d'obtenir exactement 12 piles d'environ $0.12$ (exactement $0.1201$).

Utiliser la courbe normale\index{approximation normale de la loi binomiale}. Le nombre de piles a pour espérance 10 et pour variance $20/4 = 5$, donc un écart type de $\sqrt 5 \approx 2.24$. La hauteur du sommet d'une densité normale vaut $0.4$ divisé par l'écart type, ce qui donne $0.4/2.24 \approx 0.18$. La probabilité en la moyenne est donc proche de cette valeur.

Pour 12 piles, on s'écarte de deux de la moyenne, et on multiplie par la forme normale $e^{-d^2/(2\sigma^2)} = e^{-4/10} \approx 0.67$. On obtient $0.18 \times 0.67 \approx 0.12$.

\iqwrong{Dire un demi parce que la pièce est équilibrée. Un demi est la part de piles que vous attendez, et non la probabilité d'obtenir exactement ce nombre. Il y a 21 nombres possibles, et 10 n'est que le plus probable d'entre eux.}
\iqfollow{Et pour 100 lancers~?}{L'écart type vaut 5, donc la probabilité d'obtenir exactement 50 piles est d'environ $0.4/5 = 0.08$, et la valeur exacte est $0.0796$. Comme elle diminue lorsque le nombre de lancers augmente, proportionnellement à l'inverse de sa racine carrée, il faut quatre fois plus de lancers pour la réduire de moitié.}
\end{solution}

\begin{solution}{iq:iv1:mental-technique:38}
$\binom{50}{25} \approx 1.3 \times 10^{14}$ (exactement $1.264 \times 10^{14}$), et $\binom{100}{30} \approx 2.9 \times 10^{25}$ (exactement $2.937 \times 10^{25}$).

Pour le cas central, utiliser le sommet de la loi normale\index{approximation normale de la loi binomiale}~: la probabilité d'obtenir exactement $n$ piles en $2n$ lancers équilibrés vaut environ $1/\sqrt{\pi n}$, ce qui donne $\binom{2n}{n} \approx 4^n / \sqrt{\pi n}$. Ici $4^{25} = 2^{50} \approx 1.126 \times 10^{15}$ en utilisant $1.024^5 \approx 1.126$, et $\sqrt{25\pi} \approx 8.86$, donc l'estimation vaut $1.27 \times 10^{14}$.

Pour un choix excentré, utiliser l'entropie\index{estimation de l'entropie d'un coefficient binomial}~: avec $p = 0.3$,
\[
\binom{N}{pN} \approx \frac{2^{N H(p)}}{\sqrt{2\pi N p (1 - p)}}, \qquad H(p) = -p \log_2 p - (1 - p) \log_2 (1 - p).
\]
Ici $H(0.3) \approx 0.881$, donc le numérateur vaut $2^{88.1} \approx 3.3 \times 10^{26}$. Le dénominateur vaut $\sqrt{2\pi \times 21} \approx 11.5$. On obtient $2.9 \times 10^{25}$.

\iqwrong{Utiliser $50^{25}/25!$, qui n'est une bonne borne supérieure que lorsque la sélection est une petite partie de l'ensemble. Ici elle est trop grande d'un facteur d'environ 1500.}
\iqfollow{Et $\binom{52}{26}$~?}{$\binom{52}{26} = \binom{50}{25} \times \frac{52 \times 51}{26 \times 26}$, soit un facteur de $3.92$, donc environ $4.96 \times 10^{14}$.}
\end{solution}

\begin{solution}{iq:iv1:mental-technique:39}
Environ 226 bits, et exactement $225.58$ bits. $52!$ comporte 68 chiffres.

Le nombre de mélanges est $52!$, donc le moins de bits en moyenne est $\log_2 52!$. Utiliser la formule de Stirling\index{formule de Stirling}, $\ln n! \approx n \ln n - n + \tfrac12 \ln(2\pi n)$, en base 2~:
\[
\log_2 n! \approx n \log_2 n - \frac{n}{\ln 2} + \tfrac12 \log_2 (2\pi n).
\]
Comme $52 = 32 \times 1.625$ et $\log_2 1.625 \approx 0.70$, on obtient $\log_2 52 \approx 5.70$. Les termes valent alors $52 \times 5.70 = 296.4$, $52/0.693 = 75.0$ et $\tfrac12 \log_2 327 \approx 4.2$. Le total vaut $296.4 - 75.0 + 4.2 = 225.6$.

Pour les chiffres, on multiplie par $\log_{10} 2$~: $225.6 \times 0.301 \approx 67.9$, ce qui signifie que $52!$ vaut environ $8 \times 10^{67}$, soit 68 chiffres.

\iqwrong{Multiplier 52 par 6 bits par carte, ce qui donne 312. C'est un encodage valide mais pas le plus court~: il permet qu'une carte se répète. Supprimer ce gaspillage économise plus de 80 bits.}
\iqwhy{Une graine aléatoire de 64 bits ne peut produire que $2^{64}$ mélanges sur $2^{226}$. Si les simulations dépendent du mélange complet, un générateur à petite graine n'en atteint presque aucun.}
\end{solution}

\begin{solution}{iq:iv1:mental-technique:40}
La valeur actuelle vaut environ $\$12.4$, et l'obligation vaut environ $75.2$ pour 100 de valeur nominale.

Une annuité\index{facteur d'annuité} est une perpétuité moins une perpétuité qui commence plus tard. Une perpétuité de 1~\$ par an vaut $1/r$ aujourd'hui, et une perpétuité qui commence après 30 ans vaut ce montant à l'année 30, soit $1.07^{-30}/r$ aujourd'hui. En soustrayant,
\[
\text{VA} = \frac{1 - 1.07^{-30}}{0.07}.
\]
Pour la puissance, à 7~\%, l'argent double à peu près tous les dix ans~: $1.07^{10} = 1.967$, donc proche de 2. Trois de ces doublements donnent $1.07^{30} \approx 7.6$~; la valeur exacte est $7.61$. Donc $1/7.61 \approx 0.131$, et $(1 - 0.131)/0.07 = 0.869/0.07 \approx 12.4$. La valeur exacte est $12.41$.

L'obligation verse 5 par an pendant 30 ans et rembourse 100 à la fin. Son prix vaut donc $5 \times 12.41 + 100/7.61 = 62.05 + 13.14 = 75.2$.

\iqwrong{Utiliser la date de paiement médiane~: $30 / 1.07^{15} \approx 10.9$. L'actualisation est convexe, donc les paiements précoces gagnent plus que les paiements tardifs ne perdent, et le résultat sort trop bas.}
\iqfollow{Quel paiement annuel constant rembourse un prêt de 100 sur 30 ans à 7~\%~?}{On divise par le facteur~: $100/12.41 = 8.06$ par an.}
\iqfollow{Quelle part de la valeur d'une perpétuité représente cette annuité de 30 ans~?}{Une perpétuité vaut $1/0.07 = 14.29$, donc l'annuité vaut $1 - 1/7.61$, soit environ 87~\% de cette valeur. Les paiements après l'année 30 ne valent que 13~\%.}
\end{solution}
